%%%%%%%%%%%%%%%%%%%%%%%%%%%%%%%%%%%%%%%%%%%%%%%%%%%%%%%%%%%%%%%%%%%%%%%%%%%%%%%
% Seismica submission manuscript -- shortened main paper, version 43
% Supporting detail and diagnostic figures are in supplement43.tex.
%%%%%%%%%%%%%%%%%%%%%%%%%%%%%%%%%%%%%%%%%%%%%%%%%%%%%%%%%%%%%%%%%%%%%%%%%%%%%%%

\documentclass[]{seismica}
\usepackage{tabularx}
\usepackage{placeins}

% Color and TODO command for manuscript annotations
\usepackage{xcolor}

\newcommand{\todo}[1]{%
  \par\noindent
  \colorbox{yellow!40}{%
    \parbox{0.96\linewidth}{%
      \textbf{\textcolor{red}{TODO:}} #1
    }%
  }%
  \par
}

\newcommand{\checkme}[1]{\textcolor{red}{\textbf{[CHECK: #1]}}}

% Compile from the repository latex/ directory.
\graphicspath{
  {../data/outputs/170_generate_publication_figures/}
  {../data/outputs/165_generate_video_waveform_figure/}
  {../data/outputs/145_analyze_regional_detectability/}
  {../data/outputs/140_search_for_direct_seismic_waves/}
  {../data/outputs/130_analyze_seismic_polarization/}
  {../data/outputs/120_estimate_acoustic_energetics/}
  {../data/outputs/110_measure_named_events/}
  {../data/outputs/100_measure_event_amplitudes_and_acoustic_seismic_coupling/figures/}
  {../data/outputs/100_measure_event_amplitudes_and_acoustic_seismic_coupling/}
  {../data/outputs/090_analyze_finite_distance_propagation/}
  {../data/outputs/080_analyze_planar_array_propagation/}
  {../data/outputs/040_validate_event_catalogue_with_array_processing/figures/}
  {figures/}
  {../writing/latex/figures}
}

% These boxes keep the manuscript compilable while the two remaining main
% figures are regenerated. They should not appear in a submitted PDF.
\newcommand{\pendingfigure}[1]{%
  \fbox{\parbox[c][55mm][c]{0.92\textwidth}{\centering\textbf{Pending figure}\par#1}}}

\title{Close-Range Seismo-Acoustic Observations of the 2016 Falcon 9 AMOS-6 Failure at Cape Canaveral}
\shorttitle{Close-range observations of the 2016 Falcon 9 explosion}

% Confirm department-level present affiliations for Farrell and Mehta before
% submission.
\author[1]{Glenn Thompson\thanks{Corresponding author: \texttt{thompsong@usf.edu}; ORCID: 0000-0002-9173-0097}}
\author[2,3]{Robert G. Brown}
\author[1,4]{Alexandra Farrell}
\author[2,3]{John Kiriazes}
\author[1]{Stephen R. McNutt}
\author[1,5]{Christopher Mehta}
\affil[1]{School of Geosciences, University of South Florida, Tampa, Florida, USA}
\affil[2]{NASA Kennedy Space Center, Florida, USA}
\affil[3]{Present address: Florida Space Institute, University of Central Florida, Orlando, Florida, USA}
\affil[4]{Present address: University of Alaska Fairbanks, Fairbanks, Alaska, USA}
\affil[5]{Present address: Howard University, Washington, DC, USA}

\credit{Conceptualization}{Glenn Thompson}
\credit{Methodology}{Glenn Thompson}
\credit{Software}{Glenn Thompson}
\credit{Validation}{Glenn Thompson}
\credit{Formal Analysis}{Glenn Thompson}
\credit{Investigation}{Glenn Thompson, Robert G. Brown, Alexandra Farrell,
John Kiriazes, Stephen R. McNutt, Christopher Mehta}
\credit{Resources}{Robert G. Brown, John Kiriazes}
\credit{Data Curation}{Glenn Thompson}
\credit{Visualization}{Glenn Thompson}
\credit{Project Administration}{Glenn Thompson, Robert G. Brown}
\credit{Writing - Original draft}{Glenn Thompson}
\credit{Writing - Review \& Editing}{Glenn Thompson, Robert G. Brown,
John Kiriazes, Stephen R. McNutt}

\begin{document}

\makeseistitle{
  \begin{summary}{Abstract}
  On 1 September 2016, a Falcon~9 and its AMOS-6 payload were destroyed
  during propellant loading before a planned static-fire test at Space Launch
  Complex~40, Cape Canaveral. A broadband seismometer and three-element
  infrasound array 1.4~km away recorded the complete sequence. We reconstructed
  its chronology and measured source direction, propagation, pressure, ground
  motion, and seismic/acoustic coupling. The catalogue contains 153 accepted
  pressure transients over 26~min, distributed among four prominent activity
  phases. No precursor signal was found in the seismo-acoustic data before the
  initial upper-stage explosion. The principal explosion followed the initial
  upper-stage event by about 3.6~s and produced a positive peak pressure of
  about 1.4~kPa, a sound pressure level of about 157~dB relative to the
  standard 20-$\mu$Pa reference pressure, and a hemispherical-equivalent
  acoustic-energy estimate of $2.2\times10^9$~J. High-quality array solutions
  pointed toward the launch complex and gave an apparent velocity of about
  352~m~s$^{-1}$, consistent with meteorological predictions. Pressure and
  three-component peak ground velocity were nearly proportional across 46
  well-measured catalogue entries, with a median seismic/acoustic coupling
  coefficient of $2.3\times10^{-6}$~(m~s$^{-1}$)~Pa$^{-1}$. Frequency-dependent seismic/acoustic coupling was enhanced near 4--5 and
  20--24~Hz.
  \end{summary}
  %
  \begin{summary}{Non-technical summary}
  A seismic and low-frequency sound station approximately 1.4~km from the
  launch pad recorded the 2016 Falcon~9 accident at Cape Canaveral. The
  recordings independently show when energetic activity began, how the
  sequence developed, how large its strongest signals were, and whether the
  pressure waves were consistent with propagation from the launch complex. A
  gap-free 48-h record, spanning 37~h before and 11~h after the first signal,
  was examined on all six channels, with no precursor signal identified before
  the initial upper-stage explosion.
  The accident was not a single explosion: 153 accepted multichannel pressure
  transients occurred over approximately 26~min, with four prominent activity
  phases, although some
  closely spaced entries are different peaks within the same physical event
  complex. The strongest explosion followed the first upper-stage signal by
  about 3.6~s. Independent measurements of arrival direction and speed agree with
  propagation from the launch pad under the observed weather conditions. The
  study shows that a small continuously recording geophysical array can
  provide an empirically calibrated, independently timed forensic record when access to a
  launch site is restricted and engineering telemetry is unavailable.
  \end{summary}
}

\section{Introduction}
\label{sec:introduction}

Rocket launches generate strong atmospheric and ground-motion signals, but
catastrophic launch failures are rarely recorded at close range by calibrated
seismic and infrasound instruments. Most seismo-acoustic studies of major
technological explosions have instead used regional observations, for which
atmospheric ducting and attenuation strongly influence arrival times and
amplitudes \citep{ottemoller_2008,green_2008,ceranna_2009,smink_2019,
pilger_beirut_2021,kim_pasyanos_2022}. Rocket-focused studies demonstrate
complementary capabilities: routine launches are detectable infrasonically at
large distances \citep{pilger_rocket_2021}, dense seismic deployments can
resolve launch-generated body and surface waves
\citep{schuster_rocketquake_2025}, and regional infrasound can locate and
estimate the effective yield of a large static-fire explosion
\citep{silber_scamfer_newglenn_2026}.

On 1 September 2016, a Falcon~9 was destroyed during propellant loading in the
countdown for a planned static-fire test at Space Launch Complex~40 (SLC-40),
Cape Canaveral Air Force Station (now Cape Canaveral Space Force Station),
Florida. The AMOS-6 payload was lost and the launch complex was extensively
damaged. The only publicly available original recording is the 5-min-40-s raw
USLaunchReport video and camera audio \citep{uslaunchreport_amos6_video_2016};
an edited analysis by \citet{manley_spacex_explosion_2016} was used only for
context. 

The University of South Florida (USF) BCHH station, approximately 1.4~km from
SLC-40 on adjacent Kennedy Space Center property, fortuitously recorded the
complete sequence. The initial forensic analysis was documented in an
unpublished USF technical report prepared for SpaceX in November~2016
\citep{thompson_mcnutt_fireball_2016} and subsequently presented publicly at
the 2017 AGU Fall Meeting \citep{thompson_agu_falcon9_2017}, before the
present reanalysis. The geophysical chronology was developed independently of
the later engineering root-cause determination.

The central advantage of this dataset is the combination of continuous
recording, empirical calibration, array geometry, and proximity. It provides
subsecond timing and local pressure--ground comparisons that regional records
rarely preserve, while public imagery supplies an independently inspectable
record of selected stages. Its limits are equally clear: one compact array can
constrain direction and apparent speed, but not source range or the origin of
every later transient.

Here we use the BCHH observations to test for a detectable energetic
precursor, reconstruct the accident chronology, measure the strongest events,
evaluate source direction and acoustic propagation, and characterize local
seismic/acoustic coupling. A separate regional search tests whether the
principal explosion can be associated with publicly archived records at much
greater distances. The broader objective is to show what a modest open
geophysical deployment can establish independently, and what still requires
engineering telemetry, denser spatial sampling, or detailed atmospheric
modelling. We use
\emph{close range}, rather than the strict
acoustic term \emph{near field}: a 1.4-km wavelength at 350~m~s$^{-1}$
corresponds to approximately 0.25~Hz, below the principal analysis band.
Supporting methods, diagnostics, and sensitivity results are provided in the
Supplementary Material.

\section{Data and instrumentation}
\label{sec:data}

\subsection{BCHH array and geometry}

BCHH comprised a three-component Nanometrics Trillium Compact Posthole
broadband (>0.083 Hz) seismometer near the centre of a three-element infraBSU microphone (>0.048 Hz) array with an
aperture of approximately 30~m (Figure~\ref{fig:site_array}). All six channels
were recorded continuously at 250~Hz by a GPS-synchronized Nanometrics
Centaur digitizer. The seismometer was buried approximately 0.5~m; the infrasound sensors
were mounted ~30~cm above ground on garden stakes in shrubs to suppress wind noise.

On 6 October 2016, as Hurricane Matthew approached KSC, co-author RGB 
removed the three microphones and stakes, but the cables remained on 
the ground. The cable ends were subsequently surveyed with precise GPS on
2 November 2016. Position errors of approximately 1--2~m are therefore plausible.
Source distances ranged from 1408.3 to 1440.0~m
(Supplementary Table~S1). DD2 was used as the reference
pressure channel, as it had the best signal-to-noise ratio.

The USLaunchReport video camera was located on the roof of a building in the Fire Training Area, providing a clear view of the launch complex.
We accessed this building to verify that the view matches that captured by the USLaunchReport video footage.
\todo{Add recent photo to the supplement, and cite here}


\begin{figure*}[!tbp]
  \centering
  \includegraphics[width=0.96\textwidth]{fig01_slc40_bchh_location.pdf}
  \caption{Locations used in the study (A) Position of the BCHH seismo-acoustic array
   and the USLaunchReport camera. The launch complex was SLC-40 (yellow triangle).
  (B) Three-component seismometer and three-element infrasound
  array with an aperture of approximately 30~m.}
  \label{fig:site_array}
\end{figure*}
\FloatBarrier

\subsection{Calibration, weather, video, and regional data}

The seismic response was removed using the nominal calibration for the
Trillium Compact Posthole--Centaur combination. 1-m lift tests (Supplementary Section~S1) were subsequently performed and 
determined sensitivities of 8.10, 0.69, and
10.0~counts~Pa$^{-1}$, respectively for infraBSU sensor channels DD1, DD2, and DD3 (Supplementary Table~S1).
These empirical factors \citep{thompson_mcnutt_fireball_2016} are retained here (Supplementary Table~S1).
Results from other launches reported here are consistent with these empirical factors.
The $\pm20\%$ pressure uncertainty reflects the imprecision and repeatability of the empirical lift test procedure. 
Manual arrival uncertainty was approximately 0.01~s for the clearest impulses and
greater for weak, emergent, or overlapping signals.

The USLaunchReport footage \citep{uslaunchreport_amos6_video_2016} provided the public visual and camera-audio record.
Four proprietary videos obtained from SpaceX assisted and confirmed visual
interpretation but contained no usable audio and supplied no quantitative
result reported here. Data were sought from an infrasound array at CCAFS
operated by the Air Force Technical Applications Center (AFTAC)
but were ultimately inaccessible. Regional pressure records to 1800~km
and vertical seismic records to 600~km were obtained through EarthScope FDSN
services; IU.DWPF, 97.7~km from SLC-40, was the nearest permanent station
examined \citep{asl_gsn_iu_1988}.

Kennedy Space Center meteorological observations were examined at two spatial scales. Weather-tower measurements supplied temperature, relative humidity, and near-surface winds for the close-range SLC-40--BCHH propagation correction. To characterize the vertically extended atmosphere relevant to regional infrasound propagation, we additionally examined KSC 915-MHz and 48-MHz radar wind-profiler observations and a 12:00~UTC low-resolution AMPS radiosonde sounding. The profilers constrained winds through the lower and middle atmosphere, including 48-MHz profiles at 13:04 and 13:09~UTC bracketing the principal explosion, while the AMPS sounding supplied temperature and winds into the lower stratosphere. No atmospheric ray tracing was performed. Additional details are given in Supplementary Section~S2

\FloatBarrier

\section{Methods}
\label{sec:methods}

The analysis was implemented in Python as a reproducible sequence of Jupyter
notebooks \citep{kluyver2016jupyter}. Shared configuration and intermediate
products supplied common geometry, event times, propagation corrections, and
analysis parameters. Additional implementation and sensitivity details are in
Supplementary Sections~S1--S7.

\subsection{Waveform preparation and meteorological reduction}

Response removal used a 0.02, 0.05, 80, and 100~Hz four-corner prefilter and a
60-dB water level. A centred 1.0-s moving median was then removed independently
from each pressure channel to suppress slow baseline variation without an
additional physical-amplitude high-pass filter. Seismic traces remained in
ground velocity and pressure traces in Pa.

For the close-range acoustic travel-time correction, still-air sound speed was
calculated from the mean temperature and relative humidity described in
Section~2.2. Wind vectors from the KSC weather-tower observations were
projected onto the SLC-40--BCHH propagation direction. To characterize the wind
over the nominal 0--65-m source-height interval, the available tower wind
observations within that height range were combined as a vector mean. Because
the lowest valid wind observations were near 40~m, the lowest measured wind
vector was extended downward at a constant value to estimate the nominal
0--65-m path average; the surface layer was therefore not directly sampled.
The path-parallel wind component was added to the still-air sound speed to
obtain the effective acoustic speed used for source-to-BCHH travel-time
reduction.

The tower observations were used only for this 1.4-km close-range timing
correction. Vertically extended KSC 915-MHz and 48-MHz wind-profiler
observations and the AMPS radiosonde sounding were treated separately as
regional-propagation diagnostics and were not used to alter the close-range
source-time reduction.

Three time coordinates were kept distinct throughout the analysis: observed
receiver UTC, DD2-referenced reduced arrival time, and source time. For sensor
$i$, the DD2-referenced reduced arrival time was calculated as
\begin{equation}
 t_{i,\mathrm{red}}=t_i-\frac{r_i-r_{\mathrm{DD2}}}{c_{\mathrm{eff}}},
 \label{eq:reduced_time}
\end{equation}
where $t_i$ is the observed receiver time, $r_i$ is the source-to-sensor
distance, $r_{\mathrm{DD2}}$ is the source-to-DD2 distance, and
$c_{\mathrm{eff}}$ is the effective acoustic speed derived from the tower
meteorology. Catalogue times used for event association therefore represent
arrivals reduced to DD2 rather than inferred source times. Source times were
calculated separately by subtracting the complete source-to-DD2 acoustic
travel time.

\subsection{Catalogue construction and quality control}

A gap-free 48-h interval extending from 37~h before to 11~h after the first
upper-stage signal was manually scanned using
Antelope's \texttt{dbpick} program (version~5.6) \citep{brtt_antelope}. All three pressure and all three seismic
channels were viewed together. Conspicuous impulsive signals readily
visible across the microphone array were manually picked resulting in a CSS3.0 detection table of candidate arrival times.
For each candidate event there were usually three candidate arrival times - one per microphone. The rest of the pipeline was automated. 

For association, each candidate arrival time was first reduced to the DD2 reference by
subtracting the predicted differential travel time from SLC-40 to that sensor
using the 351~m~s$^{-1}$ effective acoustic speed. A moving 0.04~s window (the smallest coincidence
window on a stable association-count plateau \todo{ADD FIGURE REFERENCE} was then applied to identify coincident 
arrivals across the array and group them into events. This simple time-based event association was performed 
using GISMO \citep{thompson_gismo_2017}. 

\todo{HOW WAS THIS PART DONE?} Candidate waveforms were subsequently reviewed for signal-to-noise ratio (SNR),
correlation, \checkme{WHAT THRESHOLDS? they seem to be given below} and a coherent counterpart on the remaining channel. The
catalogue records accepted multichannel pressure transients, not independently
verified physical explosions.

Catalogue membership was separated from waveform quality. DD1 was
sporadically noisier and, overall, correlated less consistently with the other
two pressure channels; DD2 and DD3 therefore provided the most stable
waveform-quality pair. 

\todo{MOVE TO RESULTS}The usable-waveform subset required robust peak
SNR $\geq3$ and DD2--DD3 correlation $\geq0.5$; the core subset required
SNR $\geq5$ and correlation $\geq0.7$. Time-domain propagation quality was classified independently of
these waveform criteria. The high-quality planar class required maximum
coherence $\geq0.90$, back-azimuth standard deviation
$\leq0.5^\circ$, and apparent-speed standard deviation
$\leq3.0$~m~s$^{-1}$ across nine perturbations of the event scoring window.
Intermediate solutions used corresponding thresholds of 0.75,
1.0$^\circ$, and 8.0~m~s$^{-1}$.

Independent frequency-domain Bartlett beamforming was performed using InfraPy
version~0.4.1.0, developed at Los Alamos National Laboratory (LANL)
\citep{webster_infrapy_2025}, with all three pressure channels.
The primary validation used 0.5-s windows centred on the catalogue
times and a 2--20-Hz band. If $C$ denotes the multichannel coherence for
$N=3$ sensors, the Bartlett detection statistic was
\begin{equation}
 F=\frac{(N-1)C}{1-C}.
\end{equation}
This statistic is a monotonic transformation of coherence: $F$ approaches
zero for incoherent signals and increases as the channels become more
consistent with a coherent plane wave. It therefore measures array coherence,
not pressure amplitude, acoustic energy, or source yield.

\todo{153 is a result. move it}. 
An empirical screening threshold was defined from 153 nonoverlapping
0.5-s pre-event background windows processed identically to the event-centred
windows. The threshold was taken as the 99.5th percentile of their $F$
values, i.e., the extreme upper tail of the observed background distribution.
It was used as an empirical screening criterion rather than interpreted as a
formally calibrated or scan-wide false-alarm probability.

Event-centred $F$ values provided an independent array-coherence check on
catalogue entries. A separate scan of 17,839 overlapping windows, positioned
without reference to catalogue times, was used to test catalogue recovery and
to search for additional coherent activity.

\todo{MOVE TO RESULTS}
The empirical 99.5th-percentile threshold derived from the pre-event
background windows was $F=3.02$.

\subsection{Propagation, amplitude, and seismic/acoustic coupling analyses}

A time-domain planar-wave grid search estimated back azimuth and apparent
velocity by maximizing multichannel waveform agreement. All catalogue entries
were processed, while physical interpretation emphasized quality-controlled
solutions. A fixed-source circular-wavefront calculation tested sensitivity to
the 1.4-km source range; because the three-element array provides only two
independent relative delays, it was not expected to resolve both source range
and propagation speed.

Catalogue and regression pressures were peak-to-peak amplitudes of the
differentially aligned median waveform; named-event pressures were medians
of independently measured channel extrema. Three-component vector peak ground
velocity (PGV) was measured over a pressure-defined seismic window. A log--log model
\begin{equation}
 \log_{10}p_{\mathrm{p2p}}=a+b\log_{10}(\mathrm{PGV})
 \label{eq:power_law}
\end{equation}
was fitted to entries with pressure-stack and \checkme{did we apply SNR filter again? and does some of this need to go to results?}
vector-PGV SNR both $\geq5$
and event-window PGV at least 70\% of the maximum in the saved approximately
1.1-s segment. This criterion does not establish full event/coda capture. Conditional percentile ranges were obtained by
resampling overlapping measurement groups together, with catalogue-entry
resampling retained as a comparison. Groups were defined from overlapping
sampled pressure--PGV windows; separate groups can still share coda and
calibration or propagation errors.

Frequency-dependent coupling was described by the H1
pressure-to-vertical-velocity transfer-function estimator
\begin{equation}
 H_1(f)=\frac{S_{pv}(f)}{S_{pp}(f)},
 \label{eq:h1_transfer}
\end{equation}
where $S_{pv}(f)$ is the cross-spectral density between pressure $p$ and
vertical ground velocity $v$, and $S_{pp}(f)$ is the pressure autospectral
density. The H1 estimator treats pressure as the input and vertical velocity
as the output, so $|H_1|$ has units of m~s$^{-1}$~Pa$^{-1}$ and represents
the frequency-dependent seismic/acoustic coupling coefficient. Magnitude-squared
coherence was calculated alongside $H_1$ to identify frequencies at which
pressure and vertical motion were related consistently and approximately
linearly.

Continuous estimates used the opening activity interval
(Phase~I, 13:07:11.91--13:08:21.91~UTC; see
Section~\ref{sec:temporal_evolution}). Event and continuous estimates overlap
in time and were treated as complementary rather than independent.


For the four video-interpreted opening signals later listed in
Table~\ref{tab:principal_events}---the upper-stage event, principal explosion,
payload impact, and payload explosion---acoustic exposure
was integrated over manually reviewed windows. For a progressive acoustic wave,
intensity is proportional to $p^2/(\rho c)$; assuming hemispherical spreading
and no atmospheric attenuation therefore gives
\begin{equation}
 E_{2\pi}=\frac{2\pi r^2}{\rho c}\int p^2(t)\,dt,
 \label{eq:acoustic_energy}
\end{equation}
following standard acoustic-energy treatments
\citep[e.g.,][]{johnson_aster_2005}.
The model does not explicitly correct receiver ground reflection, source
directivity, or local obstacles; changing $2\pi$ to $4\pi$ does not resolve
those effects.

Acoustic TNT equivalence is
$E_{2\pi}/(4.184\times10^6~\mathrm{J\,kg^{-1}})$,
an energy-unit conversion rather than a blast-yield estimate.
Converting radiated acoustic energy to total explosive yield additionally
requires the acoustic efficiency---the fraction of source energy radiated into
the atmosphere as sound---which depends on source physics, confinement, and
frequency content and is not constrained by this single close-range station.
Regional infrasound yield studies instead commonly rely on empirically or
physically calibrated relationships between observed signals and explosions of
known yield
\citep[e.g.,][]{ottemoller_2008,ceranna_2009,pilger_beirut_2021,kim_pasyanos_2022,
silber_scamfer_newglenn_2026}.
No confirmed regional Falcon~9 branch with a stable dominant period was
available here. Overlapping catalogue windows were not summed as independent
events. Efficiency-dependent total-energy scenarios are therefore confined to
the Supplementary Material.

\subsection{Video alignment and regional screening}

The USLaunchReport recording lacks absolute UTC timing. Its first visible upper-stage
flash was mapped to the reduced time of the first impulsive event seen across the BCHH array. Camera audio
was advanced by its 12.16-s path travel time, calculated from a 4.255-km range
and a camera-path effective speed of 349.8~m~s$^{-1}$. BCHH traces were reduced
using sensor-specific distances and 351~m~s$^{-1}$. Video supplied physical
labels for selected stages but was not used to construct the general catalogue.

Regional screening used broad physically plausible acoustic-celerity and
seismic-velocity windows. Automatic metrics compared candidate windows with matched pre- and post-event
controls and produced a permissive review list. \checkme{is this a result?}All 93 flagged phase windows
were subsequently inspected visually. 

\section{Results}
\label{sec:results}

\subsection{Atmosphere, catalogue, and temporal evolution}
\label{sec:temporal_evolution}

Seven 5-min weather-tower epochs from 13:05 to 13:35~UTC gave a mean
temperature of 26.4$\pm0.3\,^{\circ}$C and mean relative humidity of
83.9$\pm3.7\%$. The corresponding tower winds varied little among the seven
epochs, with an overall vector mean of 5.4~m~s$^{-1}$ from 144$^\circ$
(5.3--5.5~m~s$^{-1}$, 142--149$^\circ$).

For the close-range SLC-40--BCHH travel-time correction, the available tower
winds were additionally averaged over the nominal 0--65-m source-height
interval, as described in Section~2.2. This gave a path-representative
vector-mean wind of 4.6~m~s$^{-1}$ from 145$^\circ$, corresponding to a
2.6~m~s$^{-1}$ tailwind component along the propagation path. Together with
the meteorologically derived still-air sound speed of 348.3~m~s$^{-1}$, this
gave an effective acoustic speed of 351.0~m~s$^{-1}$. The corresponding travel
time over the 1419.9-m source-to-seismometer distance was approximately
4.05~s, about 0.03~s shorter than the still-air estimate.

\todo{DO A BETTER JOB OF EXPLAINING THE 138 AND 107, AND THE 51, 45, AND 57 PLANAR CLASSES. AND DO I HAVE A TABLE?}
The final catalogue contains 153 accepted multichannel pressure transient events.
These occur in four phases containing 39, 4, 90, and 6 entries,
respectively; 14 entries fall outside those intervals
(Figure~\ref{fig:overview}). The phases are descriptive intervals marking
visually separated activity clusters: Phase~I, 13:07:11.91--13:08:21.91;
Phase~II, 13:14:05--13:15:55; Phase~III, 13:19:02--13:25:00; and Phase~IV,
13:29:39--13:34:54~UTC. They are not inferred source regimes. Of the 153
entries, 138 met usable-waveform and 107 met core
criteria. The independently defined planar classes contained 51 high-quality,
45 intermediate, and 57 weak or broad solutions. Median robust peak SNR was
7.86, median DD2--DD3 correlation 0.876, and median corrected span among
associated picks 4.87~ms. These nested counts serve different analytical
purposes and should not be interpreted as counts of independent explosions.

SpaceX were particularly attentive to any precursor signals. No precursor signal 
was identified in the seismo-acoustic data during the
37~h before the initial upper-stage explosion. The visual screen was sensitive
to signals comparable to the accepted array-wide events
under the observed background; it cannot exclude weak, spatially local,
non-impulsive, internal, or otherwise poorly coupled processes. No additional
explosion-related impulse was found in the remainder of the 48-h review window,
which extends approximately 10.5~h beyond Phase~IV. A prolonged signal lasting
approximately 3.2~min followed Phase~I; it is visible on DD2 pressure as well
as the seismic channels, and from the USSpaceLaunchReport footage it coincides with
a vigorous fire across SLC-40.

\begin{figure*}[!tbp]
  \centering
  \includegraphics[width=0.98\textwidth]{../data/outputs/170_generate_publication_figures/fig03_bchh_30min_overview_zp_true.pdf}
  \caption{Thirty-minute six-channel BCHH raw data overview (no time reduction). The upper
  traces are pressure and the lower traces are three-component ground
  velocity. All traces use a
  zero-phase four-pole 0.5--25-Hz display filter and are independently scaled
  by their robust 99.3rd-percentile amplitude; amplitudes therefore cannot be
  compared between rows. Phase intervals and the prolonged post-Phase-I signal are marked. The latter
  is visible on DD2 pressure as well as on the seismic channels; the label is
  descriptive rather than a source or process classification. The shaded Phase-I--IV intervals contain 39, 4, 90, and 6
  catalogue transients, respectively; 14 transients fall outside those intervals.}
  \label{fig:overview}
\end{figure*}

The phase structure is also evident in continuous energetic
diagnostics derived from the same records (Figure~\ref{fig:sequence_energy_rates}). Five-second increments
of high-pass-filtered DD2 acoustic energy emphasize the principal explosion
and later bursts, whereas the three-component ground-motion energy proxy
separates Phases~II--IV particularly clearly. The proxy is the band-limited
increment of $\int(v_Z^2+v_R^2+v_T^2)\,dt$ and is not interpreted as radiated
seismic energy. These diagnostics support the descriptive phase
boundaries and show appreciable signal outside the short catalogue-event windows.

\begin{figure*}[!tbp]
  \centering
  \IfFileExists{../data/outputs/120_estimate_acoustic_energetics/full_sequence_cumulative_acoustic_and_ground_motion.pdf}{%
    \includegraphics[width=0.98\textwidth]{../data/outputs/120_estimate_acoustic_energetics/full_sequence_cumulative_acoustic_and_ground_motion.pdf}%
  }{%
    \pendingfigure{Continuous sequence acoustic-energy and ground-motion-proxy
    rates from Notebook 124.}%
  }
  \caption{Continuous energetic evolution of the 26-min accident sequence in
  5-s bins. (a) Hemispherical acoustic-energy increments from DD2 after a
  zero-phase 0.5-Hz high-pass filter. DD2 is shown alone because it has the
  cleanest continuous pressure record; named-event pressure, SPL, and energy
  values elsewhere in the paper retain the broader approximately
  0.05--80-Hz processed passband. (b) Three-component ground-motion energy
  proxy, the increment of $\int(v_Z^2+v_R^2+v_T^2)\,dt$ after 0.5--20-Hz
  filtering. It is a local measure of recorded ground motion, not radiated
  seismic energy. Coloured shading marks the four descriptive phase windows,
  and hatching marks the prolonged post-Phase-I signal. The especially clear
  ground-motion bursts support the adopted Phase~II--IV
  boundaries.}
  \label{fig:sequence_energy_rates}
\end{figure*}
\FloatBarrier

\subsection{Opening sequence and key events}

Four opening-sequence signals at the start of phase~1 were assigned physical interpretations from the
mapped video chronology (Table~\ref{tab:principal_events}). The nominal
upper-stage source time was 13:07:11.91~UTC. The principal explosion followed about 3.6~s
later and produced by far the largest acoustic and seismic amplitudes
(Table~\ref{tab:principal_events}). The two payload-related signals were separated
by 0.568~s. Event~024 at nominal source time 13:07:26.46~UTC had comparable
amplitude but no unique physical process could be assigned from the available video; it remains a
catalogue event rather than a fifth named stage.

Figure~\ref{fig:multimodal} shows that the first visible flash, reduced camera
audio, pressure, and ground motion are mutually consistent within the timing
limitations. BCHH source times depend on the 351~m~s$^{-1}$ reduction, and the
video's mapped UTC remains conditional on the visual anchor. Absolute source
times are therefore reported to 0.01~s as nominal values; relative intervals
between consistently reduced picks are better constrained.
Figure~\ref{fig:named_event_video_pressure} complements that continuous
alignment by comparing the pressure signatures of the four physically
interpreted signals with selected visual states.

\begin{figure*}[!tbp]
  \centering
  \includegraphics[width=0.98\textwidth]{../data/outputs/170_generate_publication_figures/fig04_multimodal_alignment_upper_stage.pdf}
  \caption{Video, camera audio, pressure, and ground-velocity alignment for
  the upper-stage event. Camera audio was advanced by its path travel time at 348.8~m~s$^{-1}$and
  BCHH traces by sensor-specific travel times at 351~m~s$^{-1}$. Pressure is baseline corrected;
  seismic traces are response corrected without an additional display filter.
  Traces are independently normalized. Video frames copyright 2016
  USLaunchReport.com/Michael Wagner; reproduced with permission. The camera audio is ~0.005~s late suggesting the wind may be ~0.15~m~s$^{-1}$ 
  slower than assumed.
  \citep{uslaunchreport_amos6_video_2016}.}
  \label{fig:multimodal}
\end{figure*}

\begin{figure*}[!tbp]
  \centering
  \includegraphics[width=0.88\textwidth]{../data/outputs/165_generate_video_waveform_figure/165_video_named_event_waveforms.pdf}
  \caption{Selected visual states and event-centred pressure waveforms of the
  opening failure sequence. Each composite panel pairs a manually selected
  representative video frame with detrended microphone data reduced at 351~m~s$^{-1}$ for
  (a) the upper-stage failure, (b) the principal explosion, (c) payload
  descent and payload impact, and (d) the payload explosion. Frames were
  selected for visual clarity and do not, in general, represent the inferred
  physical origin times. DD2 (solid blue) is emphasized because it provides
  the highest-quality pressure record; DD1 (orange dotted) and DD3 (green
  dash-dotted) are retained for intersensor comparison. For each trace, the
  vertical dotted line at zero is the predicted channel-specific acoustic
  arrival calculated using 351~m~s$^{-1}$. Event time is relative to the upper-stage
  source time. These sensor-relative display windows and
  local baselines follow Supplementary Section~S5 and need not reproduce the
  unshifted named-event measurement windows. Video frames copyright 2016 USLaunchReport.com/Michael Wagner; reproduced with
  permission \citep{uslaunchreport_amos6_video_2016}.}
  The principal explosion (b) resembles a Friedlander waveform, produced by a blast.
  It arrives approximately 0.12~s earlier than the prediced arrival time, suggesting an average 
  propagation speed of 362~m~s$^{-1}$, rather than 351~m~s$^{-1}$. Similar, the payload impact (c) and
  explosion (d) arrive earlier than predicted, indicating travel speeds of 353-354 m~s$^{-1}$. 
  \label{fig:named_event_video_pressure}
\end{figure*}

\begin{table*}[!tbp]
\caption{Physically interpreted opening-sequence signals. Source times are
nominal values rounded to 0.01~s; their absolute accuracy is limited by the
acoustic travel-time reduction. DD2 arrival times are the corresponding
pressure-arrival times at the DD2 infrasound microphone, calculated using the
1419.9-m source-to-DD2 distance and the meteorologically derived effective
acoustic speed of 351.0~m~s$^{-1}$. Pressure values are medians of three
sensor-specific measurements. Peak level is referenced to 20~$\mu$Pa.
Event~024 and all catalogue-quality measurements are reported in
Supplementary Section~S5. PGVs trace to catalogue rows 2, 13, 22 and 23,
respectively; these are representative window measurements, and unique
initial/principal pulse assignments remain unverified
(Supplementary Section~S5).}
\label{tab:principal_events}
\centering
\small
\begin{tabularx}{\textwidth}{Xccccccc}
\hline
Signal &
\shortstack{Source time\\(UTC)} &
\shortstack{DD2 arrival\\time (UTC)} &
$p_+$ (Pa) &
$p_-$ (Pa) &
$p_{\mathrm{p2p}}$ (Pa) &
\shortstack{$L_{\mathrm{peak}}$\\(dB)} &
\shortstack{Vector PGV\\(m s$^{-1}$)} \\
\hline
Upper-stage event &
13:07:11.91 &
13:07:15.96 &
38.8 &
$-25.4$ &
77.1 &
125.8 &
$7.85\times10^{-5}$ \\

Principal explosion &
13:07:15.51 &
13:07:19.56 &
1390 &
$-160$ &
1550 &
156.8 &
$3.57\times10^{-3}$ \\

Payload impact &
13:07:24.42 &
13:07:28.47 &
279.1 &
$-115.3$ &
390.7 &
142.9 &
$6.98\times10^{-4}$ \\

Payload-associated explosion &
13:07:24.99 &
13:07:29.03 &
297.0 &
$-113.4$ &
410.6 &
143.4 &
$9.02\times10^{-4}$ \\
\hline
\end{tabularx}
\end{table*}
\FloatBarrier

\subsection{Waveforms, array propagation, and seismic/acoustic coupling}

The catalogue was divided into 51 high-quality planar, 45 intermediate, and
57 weak/broad events according to the multichannel pressure coherence and
planarity criteria described above. The 51 high-quality planar events show
broadly repeatable impulsive compression followed by a smaller and longer
rarefaction, whereas vertical ground velocity has a substantially longer
oscillatory coda (Figure~\ref{fig:waveform_ensemble}). The strongest pressure
pulses are Friedlander-like in form; classical symmetric N-wave morphology
was not observed. Many weaker pressure impulses approached the seismic and
acoustic noise floors.

\begin{figure*}[!tbp]
  \centering
  \includegraphics[width=0.88\textwidth]{../data/outputs/100_measure_event_amplitudes_and_acoustic_seismic_coupling/figures/high_quality_normalized_pressure_and_dhz_waveforms.pdf}
  \caption{Independently normalized pressure and vertical-ground-velocity
  waveforms for the 51 events in the high-quality planar class. Pale curves
  are individual events, the dark curve is the pointwise median, and shading
  spans the 16th--84th percentiles.}
  \label{fig:waveform_ensemble}
\end{figure*}

Two different propagation-speed measures are used below. Source-to-array
average speed is inferred from travel time over the approximately 1.42-km
SLC-40--BCHH path, whereas array-crossing apparent velocity is measured
locally across the approximately 30-m BCHH aperture. The latter is also
sensitive to metre-scale uncertainty in reconstructed sensor positions.
These quantities are therefore reported separately.

In the event-centred InfraPy Bartlett analysis, median back azimuth was
200.7$^\circ$ compared with the geometric 199.4$^\circ$; 145 of 153 events
were within 5$^\circ$ of SLC-40. Median apparent velocity was
354.6~m~s$^{-1}$, and 113 events exceeded the empirical background threshold
$F=3.02$. The blind scan matched 129 catalogue entries (84.3\%) one-to-one;
recovery increased to 75 of 77 entries with event-centred $F\geq5$ and
41 of 42 with $F\geq10$. Non-recovery was concentrated among weak or
overlapping signals, and local maxima from overlapping scan windows were not
interpreted as additional explosions.

The independent time-domain analysis gave a high-quality
coherence-weighted direction of 200.7$^\circ$ from projected grid north and
an apparent velocity of 352.0~m~s$^{-1}$, within approximately
1~m~s$^{-1}$ of the meteorologically predicted effective speed. The
geographic-to-grid bearing conversion is approximately 0.20$^\circ$; small
direction offsets also depend on the adopted reference point. Conditional
sensor-position tests show that 1--2-m reconstruction errors across the
approximately 30-m aperture can exceed the observed direction and velocity
dispersions (Supplementary Section~S4). The finite-distance calculation did
not resolve range and slightly reduced coherence; the planar result is
therefore preferred.

For the principal explosion, the preferred time-domain array-crossing
solution was 352.5~m~s$^{-1}$ with coherence 0.9984. This is only
0.5~m~s$^{-1}$ above the coherence-weighted high-quality value and
approximately 0.13 of the 3.85~m~s$^{-1}$ high-quality velocity dispersion.
A separate short-window Bartlett solution gave a higher apparent velocity of
approximately 369.4~m~s$^{-1}$, with $F\simeq4.49$.

The source-to-array timing gives a different result. Relative to the adopted
351~m~s$^{-1}$ source-to-DD2 reduction, the principal explosion arrived
approximately 0.12~s earlier than predicted. Over the 1419.9-m
source-to-DD2 path, this residual corresponds to a path-averaged propagation
speed of approximately 362~m~s$^{-1}$, conditional on the adopted source-time
reference. Thus the principal event has a faster source-to-array average speed
than its preferred local array-crossing velocity.

\begin{figure*}[!tbp]
  \centering
  \IfFileExists{../data/outputs/080_analyze_planar_array_propagation/combined_04b_08_array_summary.pdf}{%
    \includegraphics[width=0.98\textwidth]{../data/outputs/080_analyze_planar_array_propagation/combined_04b_08_array_summary.pdf}%
  }{%
    \pendingfigure{Combined time-domain planar and event-centred Bartlett
    array summary through source UTC. Re-run Notebook 080 to generate this
    product.}%
  }
  \caption{Array-processing results through the 153-transient sequence.
  (a) Time-domain planar-wave back azimuth and (b) apparent velocity,
  classified using coherence and solution stability. (c) Independent
  event-centred Bartlett $F$ statistic. Blue circles, orange triangles, and
  grey points denote high-quality, intermediate, and weak or broad time-domain
  solutions, respectively. Dashed lines show the geometric SLC-40 direction,
  the meteorologically predicted effective acoustic speed, and the empirical
  99.5th-percentile Bartlett background threshold, respectively. In panel (c),
  larger $F$ indicates a more coherent plane-wave fit across the pressure
  array; $F$ is not a measure of pressure amplitude or explosion energy.
  Open squares identify time-domain velocity-grid-edge solutions, and shading
  marks the four activity phases. Planar azimuths use projected grid north;
  the configured seismometer-reference geometric bearing is converted to grid
  north. Times have been reduced to estimated source UTC.}
  \label{fig:array_summary}
\end{figure*}

The measured positive overpressure provides an independent weak-shock
consistency calculation. If the principal-event median positive pressure of
approximately 1.39~kPa is interpreted as a normal-shock pressure jump, the
Rankine--Hugoniot relation gives $M=1.0059$. With the
348.33~m~s$^{-1}$ still-air sound speed this corresponds to a local shock
speed of 350.37~m~s$^{-1}$, or 353.02~m~s$^{-1}$ after adding the
2.65~m~s$^{-1}$ path-parallel tailwind. By contrast, a local shock speed of
approximately 369.4~m~s$^{-1}$ would require an overpressure of order
13~kPa, roughly an order of magnitude larger than measured.

The payload-impact and payload-explosion signals had positive pressures of
order 0.28--0.30~kPa. If interpreted in the same normal-shock framework,
these correspond to $M\approx1.001$ and to nonlinear speed excesses of less
than approximately 1~m~s$^{-1}$ at BCHH. Such differences are smaller than
the practical velocity resolution imposed by sensor-position and
meteorological uncertainties. Supplementary Section~S5 additionally examines
the principal-event overpressure under a deliberately simple inward
$\Delta p\propto r^{-1}$ sensitivity model.

Forty-six catalogue entries occupying 45 overlapping-measurement groups met
the joint criteria for pressure--PGV regression. Measurement groups were
defined by touching or overlapping sampled measurement windows, including
connections through unselected bridging rows. The fitted slope was 0.962
(overlap-group bootstrap 2.5--97.5 percentiles 0.871--1.011),
$R^2=0.953$, and Spearman correlation 0.931
(Figure~\ref{fig:pressure_pgv}). Excluding all selected rows in the principal
measurement group gave $b=0.943$ and $R^2=0.905$ for 44 entries. The
near-unit scaling therefore persisted with lower explained variance. The
median selected-entry pressure-to-vector-PGV ratio was
$4.34\times10^5$~Pa per (m~s$^{-1}$), equivalently
$2.31\times10^{-6}$~(m~s$^{-1}$)~Pa$^{-1}$. These resampling ranges exclude
calibration uncertainty and dependence between nominally separate
measurement groups.

\begin{figure*}[!tbp]
  \centering
  \includegraphics[width=0.82\textwidth]{../data/outputs/100_measure_event_amplitudes_and_acoustic_seismic_coupling/figures/publication_pressure_pgv_regression.pdf}
  \caption{Peak-to-peak pressure versus three-component vector PGV for 46
  measurement-quality catalogue entries. Qualified events are shaded by
  pressure SNR; the principal-explosion star uses the same scale. The solid
  curve is the fitted power law. The dashed line shows the median proportional
  relation, $p_{\mathrm{p2p}}=4.34\times10^5\,\mathrm{PGV}$, equivalently
  $\mathrm{PGV}=2.31\times10^{-6}p_{\mathrm{p2p}}$, in SI units.}
  \label{fig:pressure_pgv}
\end{figure*}

Event-window and continuous analyses both showed enhanced vertical response
per unit pressure near 4--5 and 20--24~Hz, with a weaker 13--14-Hz feature
(Figure~\ref{fig:transfer}). The H1 pressure-to-vertical-ground-velocity
transfer-function magnitude reached approximately
$5$--$7\times10^{-6}$~(m~s$^{-1}$)~Pa$^{-1}$ in the principal coupling
bands, with elevated coherence across parts of the same frequency ranges.
The complementary event-window and continuous Phase~I estimates therefore
identify repeatable frequency-dependent seismic/acoustic coupling at BCHH.

\begin{figure*}[!tbp]
  \centering
  \includegraphics[width=0.94\textwidth]{../data/outputs/100_measure_event_amplitudes_and_acoustic_seismic_coupling/figures/publication_pressure_dhz_transfer_and_coherence.pdf}
  \caption{Empirical frequency-dependent seismic/acoustic coupling coefficient
  at BCHH. The H1 pressure-to-vertical-ground-velocity magnitude is shown for
  the high-quality event ensemble and for continuous Phase~I, together with
  magnitude-squared coherence. Agreement between the complementary
  event-window and continuous estimates identifies repeatable site-response
  bands near 4--5 and 20--24~Hz.}
  \label{fig:transfer}
\end{figure*}

\subsection{Acoustic energetics and regional search}

The continuous sequence-rate view in
Figure~\ref{fig:sequence_energy_rates} complements the event-window
measurements below by showing when acoustic energy and the local ground-motion
proxy accumulated, without treating every catalogue pick as an independent
energy release.

Within the processed infraBSU passband (approximately 0.05--80~Hz), the
manually reviewed principal-explosion window had a median positive-peak sound
pressure level of 156.8~dB re 20~$\mu$Pa and a median
hemispherical-equivalent radiated acoustic energy of
$2.2\times10^9$~J across DD1--DD3
(range $1.93$--$2.51\times10^9$~J). Expressed only as an energy-unit
conversion, this is equivalent to approximately 525~kg TNT.

In the separate de-duplicated sequence analysis, 51 selected event windows
merged into 47 nonoverlapping complexes. Their cumulative median radiated
acoustic energy was $2.27\times10^9$~J
(range $1.95$--$2.55\times10^9$~J), equivalent to 543~kg TNT
(range 467--610~kg). The cumulative value covers only the selected
measurement intervals rather than the complete 26-min sequence. A separate
source of uncertainty is the approximately $\pm20\%$ reproducibility of the
pressure calibration procedure; because acoustic energy scales with squared
pressure, this corresponds to a median-energy range of approximately
$(1.4\text{--}3.2)\times10^9$~J.

These acoustic-energy estimates can be compared with, but should not be
equated to, the chemical-energy reservoir of the launch vehicle. Published
Falcon~9 v1.2 specifications give nominal propellant masses of approximately
411~t for the first stage and 111.5~t for the second stage, for approximately
522.5~t of LOX/RP-1 at full load
\citep{<Falcon9_datasheet_key>}. Using the nominal LOX/RP-1 mixture ratio and
the heating value of RP-1 gives a full-load chemical-energy reservoir of order
$6$--$7\times10^{12}$~J. The exact propellant inventory aboard AMOS-6 at the
time of failure is not publicly documented, and the accident occurred during
the propellant-loading sequence; the full-load value is therefore an upper
bound rather than an estimate of the instantaneous onboard energy.

The principal-explosion acoustic energy is therefore only about
$3\times10^{-4}$ of the full-load chemical-energy reservoir, or approximately
0.03\%. This small fraction should not be interpreted as an anomalously low
acoustic efficiency, because it spans two distinct conversion steps:
chemical energy to effective impulsive blast energy, and blast energy to
radiated acoustic energy. For comparison, the 2026 New Glenn static-fire
explosion had an infrasound-derived effective blast yield of approximately
0.131~kt TNT ($\sim5.5\times10^{11}$~J), and comparison with plausible onboard
chemical-energy reserves implied that approximately 3--4\% of the available
chemical energy participated in the rapid pressure-producing blast under a
full-fill scenario \citep{silber_scamfer_newglenn_2026}. That quantity is a
blast-coupling fraction, not a radiated acoustic-energy efficiency.

If Falcon~9 had exhibited a comparable 3--4\% conversion from available
chemical energy to impulsive blast energy, a fully loaded vehicle would imply
an effective blast energy of approximately
$(2.0$--$2.7)\times10^{11}$~J, or about 48--64~t TNT equivalent. The measured
$2.2\times10^9$-J acoustic energy would then represent approximately
0.8--1.1\% of that effective blast energy. This is only a sensitivity
comparison, because neither the actual Falcon~9 fill fraction nor its
blast-coupling fraction is independently known; it is not a yield estimate.
The full energy-partition calculation is given in Supplementary
Section~S5.

Regional screening began with 171 station--sensor targets, of which 165
produced usable waveform records. Automatic screening of the expected arrival
windows then flagged 93 candidate phase windows at 67 stations for manual
review. The screening was intentionally permissive: for pressure data it used
a broad 180--350~m~s$^{-1}$ celerity range, divided into nominal
thermospheric-or-slow (180--260~m~s$^{-1}$), stratospheric
(260--310~m~s$^{-1}$), and tropospheric (310--350~m~s$^{-1}$) bins. These
labels define screening windows only and do not identify actual atmospheric
propagation paths. Multiple phase windows could therefore be flagged at a
single station, and the 93 flags should be regarded as review candidates
rather than detections.

All 93 flagged windows, including IU.DWPF, were visually reviewed. Three
isolated 0.2--1.0-Hz pressure packets at N4.257A, N4.S54A, and N4.R55A
(381--1079~km) had apparent celerities within the broad infrasound screening
range and peak amplitudes of 0.162--0.352~Pa. However, no coherent regional
moveout was identified in the corresponding 98-station northern record
section, and none of the three packets was accepted as a Falcon~9 detection.
They were therefore not used for period--yield analysis. Supporting maps,
record sections, screening criteria, and review results are in Supplementary
Section~S6.

The upper-air observations showed strongly height-dependent winds, with
southerly to southwesterly flow through much of the lower troposphere, stronger
westerly flow near 10--15~km altitude, and a reversal to easterly flow in the
lower stratosphere. The AMPS temperature profile showed decreasing sound speed
through the troposphere followed by recovery above the tropopause, producing a
strongly height- and azimuth-dependent effective sound-speed structure. These
conditions are consistent with substantial directional refraction and possible
regional acoustic shadowing, but no atmospheric propagation model was used to
predict specific return ranges.

\FloatBarrier

\section{Discussion}
\label{sec:discussion}

\subsection{Forensic chronology}

The BCHH record supplies a chronology independent of vehicle telemetry. One
question posed during the initial forensic review was whether the record
contained an external precursory acoustic or seismic signal that might indicate
an externally initiated event. The 48-h six-channel review found no such signal
before the initial upper-stage explosion, consistent with the November~2016
forensic report \citep{thompson_mcnutt_fireball_2016} and the later 2017 AGU
presentation \citep{thompson_agu_falcon9_2017}. This result does not exclude
weak, internal, non-impulsive, or poorly coupled processes, but provides no
geophysical evidence for a conspicuous external precursor.

Source-reduced arrivals and the video support the sequence of initial
upper-stage event, principal explosion, payload impact, and later payload
explosion. The video is anchored to the initial event, so it distinguishes
stages within the closely spaced catalogue without independently fixing UTC.
SpaceX later attributed the initiating failure to a composite overwrapped
pressure vessel in the second-stage liquid-oxygen tank
\citep{spacex_anomaly_2017,nasa_sma_falcon9}. BCHH cannot identify that
internal mechanism. Its abrupt onset is compatible with, but does not prove,
the reported 93-ms transition from anomalous telemetry to second-stage loss.

\subsection{Acoustic propagation and weak-shock interpretation}

Several independent observations are consistent with nonlinear blast-wave
propagation during the principal explosion. The pressure waveform shows rapid
compression followed by a longer rarefaction and is Friedlander-like rather
than a classical symmetric N-wave. More importantly, the principal arrival
reaches BCHH approximately 0.12~s earlier than predicted by the adopted
351~m~s$^{-1}$ source-to-array acoustic reduction. Over the approximately
1419.9-m source-to-DD2 path, this timing difference corresponds to a
path-averaged propagation speed of approximately 362~m~s$^{-1}$. Although
the exact residual also depends on source-time assignment and meteorological
uncertainty, its sign and magnitude are consistent with propagation that was
faster than the nominal effective acoustic speed over at least part of the
source--receiver path.

This path-averaged speed should not be confused with the apparent velocity
measured locally across the approximately 30-m BCHH array. The preferred
time-domain estimate for the principal explosion is 352.5~m~s$^{-1}$, close
to the meteorologically predicted effective acoustic speed at the receiver.
That result does not argue against initially supersonic propagation: a
nonlinear blast wave is expected to decelerate as its overpressure decreases,
so a disturbance that travelled appreciably faster close to the vehicle could
have approached the ambient acoustic speed by 1.4~km range. The higher
short-window Bartlett estimate of approximately 369.4~m~s$^{-1}$ is in the
same sense as the early source-to-array arrival, although its modest
$F\simeq4.49$ and sensitivity to array geometry make it less reliable as a
quantitative local-speed estimate.

The measured overpressure provides an independent constraint on the wave state
at BCHH. If the principal-event median positive pressure of approximately
1.39~kPa is interpreted as a normal-shock pressure jump, it corresponds to
$M\approx1.006$, i.e. only a very weak shock at the receiver. Thus the
pressure measurement is consistent with a wave that remained slightly
nonlinear at BCHH while having decayed substantially from a stronger
near-source blast. Conversely, interpreting the approximately
369~m~s$^{-1}$ Bartlett value literally as the local shock speed would require
an overpressure of order 13~kPa, about an order of magnitude larger than
observed, so that value is unlikely to represent the true local shock speed.

The simple inward sensitivity calculation in Supplementary Section~S5
supports the same interpretation. Extending the measured overpressure inward
with $\Delta p\propto r^{-1}$ produces substantially larger near-source
overpressure and a path-averaged speed approaching 360~m~s$^{-1}$ from
10~m to BCHH. This is not a source-pressure inversion, because the spreading
law is not validated in the near field, but it demonstrates that appreciably
nonlinear propagation close to the vehicle is compatible with a nearly
acoustic wavefront at the receiver.

Taken together, the early arrival, Friedlander-like waveform, measured
overpressure, array propagation estimates, and weak-shock sensitivity
calculation favor an interpretation in which the principal explosion generated
a nonlinear blast wave that propagated supersonically close to the vehicle and
decayed toward ordinary acoustic propagation before reaching BCHH. The data do
not require a strongly supersonic wavefront at 1.4~km, but neither does the
near-acoustic local array speed provide evidence against supersonic propagation
earlier along the path. We therefore retain 351~m~s$^{-1}$ as the reference
speed for catalogue source-time reduction while treating the approximately
0.12-s principal-event residual as evidence that this reduction does not fully
represent the propagation of the strongest blast.

\subsection{Seismic/acoustic coupling and acoustic energy}

One of the more transferable results is the empirical local
seismic/acoustic coupling coefficient. Across 46 measurement-quality entries,
peak-to-peak pressure scaled nearly linearly with three-component vector PGV:
the fitted exponent was 0.962, and the median coupling relation was
\[
 p_{\mathrm{p2p}}=4.34\times10^5\,\mathrm{PGV},
 \qquad
 \mathrm{PGV}=2.31\times10^{-6}p_{\mathrm{p2p}},
\]
with pressure in Pa and velocity in m~s$^{-1}$. Removing the principal
measurement group changed the exponent only to 0.943. The scaling therefore
does not arise solely from the largest event, although it remains an empirical
BCHH relation rather than a universal pressure-to-ground-motion conversion.
This velocity-to-pressure ratio is directly comparable in form to
seismic/acoustic coupling coefficients derived from Hunga Tonga--Hunga Ha'apai
pressure waves
\citep{anthony_hunga_coupling_2023,macpherson_hunga_coupling_2025}, while
remaining conceptually distinct from energy-partition metrics such as the
volcano acoustic-to-seismic ratio (VASR) \citep{johnson_aster_2005}.

The H1 analysis reveals frequency dependence hidden by the broadband
peak-amplitude relation. Vertical response per unit pressure was strongest near
4--5 and 20--24~Hz, with a weaker 13--14-Hz feature; the principal-band H1
magnitude was about
$5$--$7\times10^{-6}$~(m~s$^{-1}$)~Pa$^{-1}$. It need not equal the
$2.31\times10^{-6}$ broadband coupling coefficient because H1 uses vertical
motion and spectral averages within individual frequency bands. Agreement
between event-window and continuous estimates supports a repeatable local
seismic/acoustic response, comparable in scale to controlled-explosion,
volcanic, and Hunga Tonga--Hunga Ha'apai measurements
\citep{novoselov_ground_coupling_2020,anthony_hunga_coupling_2023,
macpherson_hunga_coupling_2025}. A second seismic station and a shallow
velocity model would be needed to relate this response more directly to local
material properties and subsurface structure.

The pressure calibration remains an important limitation on both the
coupling and energetic estimates, but the lift-test calibration is supported
by the independent launch comparison. For the likely 1-inch DD3 head, its
nominal sensitivity of 9.2~counts~Pa$^{-1}$ is close to the adopted
10.0~counts~Pa$^{-1}$ lift-test value: using the nominal value would raise
DD3-derived pressures by only 8.7\%. This agreement does not establish an
absolute laboratory calibration, and the anomalously low DD2 sensitivity
remains an important qualification, but it supports use of the empirical
factors within the stated $\pm20\%$ pressure uncertainty.

The principal explosion contained a median hemispherical-equivalent radiated
acoustic energy of $2.2\times10^9$~J, close to the
$2.27\times10^9$~J accumulated across all selected nonoverlapping event
complexes. Under the adopted processing and $2\pi$ radiation model, the
principal explosion therefore dominates the selected acoustic-energy budget.
These quantities are finite-passband, model-equivalent radiated acoustic
energies rather than blast yields. Absolute calibration, possible sensor
over-range response, baseline choice, receiver reflection, source directivity,
and local obstacles are correspondingly more important than the reported
numerical precision.

A useful interpretation requires separating three energy scales:
the chemical energy stored in the propellants, the fraction participating
promptly in the pressure-producing blast, and the still smaller fraction
radiated acoustically. Nominal Falcon~9 Full Thrust propellant capacity
corresponds to a full-load chemical-energy reservoir of order
$6$--$7\times10^{12}$~J. The measured principal-explosion acoustic energy is
therefore only about $3\times10^{-4}$, or 0.03\%, of that reservoir. This
chemical-to-acoustic ratio spans both conversion steps and should not be
interpreted as an acoustic efficiency. Moreover, the vehicle failed during
propellant loading, so the actual onboard chemical-energy reservoir was below,
or at most equal to, the nominal full-load value.

The distinction is illustrated by the 2026 New Glenn static-fire explosion.
Regional infrasound observations gave an effective blast yield of approximately
0.131~kt TNT, or $5.5\times10^{11}$~J, and comparison with plausible onboard
chemical-energy reserves implied a blast-coupling fraction of approximately
3--4\% for a full-fill scenario
\citep{silber_scamfer_newglenn_2026}. This is the fraction of chemical energy
represented by the inferred impulsive blast, not the fraction radiated as
acoustic energy. The New Glenn study therefore constrains a different step in
the energy partition from the close-range Falcon~9 measurement.

As a sensitivity comparison, applying the same 3--4\% blast-coupling fraction
to a fully loaded Falcon~9 would give an effective impulsive blast energy of
approximately $(2.0$--$2.7)\times10^{11}$~J, or about 48--64~t TNT
equivalent. The measured $2.2\times10^9$-J radiated acoustic energy would then
correspond to approximately 0.8--1.1\% of that effective blast energy. These
values are not an independent Falcon~9 yield estimate: neither the actual
propellant fill fraction nor the chemical-to-blast coupling is known. They do,
however, demonstrate that the apparently small 0.03\% chemical-to-acoustic
ratio does not require an exceptionally inefficient acoustic source. A
two-stage partition,
\[
 E_{\mathrm{chemical}}
 \rightarrow E_{\mathrm{blast}}
 \rightarrow E_{\mathrm{acoustic}},
\]
with a few percent entering the impulsive blast and of order one percent of
that blast energy appearing in the present radiated-acoustic estimate, is
consistent with both the Falcon~9 measurements and the New Glenn comparison.
The detailed sensitivity calculation is given in Supplementary Section~S5.

No regional counterpart to the Falcon~9 explosion was confirmed. Of 171
station--sensor targets examined, 165 produced usable waveform records; the
automatic screen subsequently flagged 93 candidate phase windows at 67
stations for manual review. These were permissive review candidates rather
than detections, and multiple windows could occur at a single station. All 93
were visually reviewed. Three isolated pressure packets at N4.257A, N4.S54A,
and N4.R55A had source-to-receiver celerities within the broad range expected
for regional atmospheric propagation (approximately 0.18--0.33~km~s$^{-1}$;
typical stratospheric arrivals are approximately 0.28--0.32~km~s$^{-1}$),
but no coherent regional moveout was identified, and none was accepted as a
Falcon~9 detection. The event therefore cannot support a regional
period--yield analysis of the kind applied to New Glenn
\citep{silber_scamfer_newglenn_2026}.

The regional non-detection is not necessarily inconsistent with the strong
close-range signal. Contemporaneous KSC wind-profiler and AMPS sounding data
show a strongly height- and azimuth-dependent atmosphere, including decreasing
temperature-dependent sound speed through the troposphere, strong westerly
flow near 10--15~km, and a wind reversal in the lower stratosphere. These
observations indicate conditions favorable for substantial directional
refraction and are consistent with possible acoustic shadowing or strongly
favored propagation corridors. Because the measured sounding extends only
through the lower stratosphere and no ray-tracing or full-wave propagation
calculation is included here, we do not use the atmospheric structure to
predict specific return ranges or explain individual station
non-detections. Wind-dependent spectral filtering provides an additional
reason not to transfer regional period--yield relations without propagation
modelling \citep{silber_wind_filtering_2026}.

\subsection{Monitoring implications and limitations}

A single stand-off site constrained onset, chronology, relative amplitude,
source direction, acoustic speed, and local seismic/acoustic coupling while pad
access was restricted. A permanent KSC seismo-acoustic network could build on
that result. Several compact pressure arrays paired with separated
three-component seismic stations would provide triangulation and path sampling,
separate direct ground waves from local seismic/acoustic conversion, and test
whether the BCHH seismic/acoustic coupling coefficient transfers between sites. Continuous
recording would preserve the background needed for both real-time situational
awareness and later forensic analysis.

The main limitations are the single site, metre-scale reconstructed sensor
positions, empirical pressure calibration, pressures beyond nominal sensor
ranges, manual event identification, uncertain absolute video timing, and
first-order radiation assumptions. The provisional $\pm20\%$ pressure
calibration uncertainty affects all absolute pressure-dependent results.
Baseline, geometry, propagation, and possible over-range nonlinearity add
uncertainties outside that percentage. Nearby-launch raw-count checks support
relative DD2/DD3 stability, but do not identify DD2's nominal head or establish
an absolute calibration. The timing, direction, and scaling results remain
useful; absolute energy, yield, and seismic-phase interpretations require more
caution.

\section{Conclusions}
\label{sec:conclusions}

The BCHH array captured an unusually complete close-range record of the 2016
Falcon~9 AMOS-6 failure. It resolved 153 accepted multichannel pressure
transients over 26~min. No precursor signal was identified during the 37~h
before the initial upper-stage explosion. The principal explosion followed
that event by about 3.6~s and produced by far the strongest pressure and ground
motion. Within the processed infraBSU passband (approximately 0.05--80~Hz),
its median positive-peak sound pressure level was 156.8~dB re 20~$\mu$Pa and
its median hemispherical-equivalent radiated acoustic energy was
$2.2\times10^9$~J.

The strongest event also shows evidence for nonlinear blast-wave propagation.
Its pressure waveform is Friedlander-like, and its arrival at BCHH is about
0.12~s earlier than predicted by the adopted 351~m~s$^{-1}$ acoustic
reduction, corresponding to a source-to-array path-average speed of
approximately 362~m~s$^{-1}$. In contrast, the preferred local
array-crossing speed at 1.4~km range is only 352.5~m~s$^{-1}$. The combination
is consistent with a blast wave that propagated supersonically closer to the
vehicle and decayed toward the ambient acoustic speed with range. The measured
1.39-kPa positive overpressure corresponds to only $M\approx1.006$ if treated
as a normal-shock jump, indicating that any shock remaining at BCHH was weak.

The measured acoustic energy is distinct from explosive yield. A fully loaded
Falcon~9 contains of order $6$--$7\times10^{12}$~J of propellant chemical
energy, so the principal-event acoustic energy is only about 0.03\% of that
reservoir. This ratio spans both chemical-to-blast coupling and
blast-to-acoustic radiation. For comparison, the 2026 New Glenn static-fire
explosion had an infrasound-derived effective blast yield of approximately
0.131~kt TNT and a reported full-fill blast-coupling fraction of about
3--4\% \citep{silber_scamfer_newglenn_2026}. Applying the same coupling
fraction to a fully loaded Falcon~9 gives an illustrative effective blast
energy of approximately 48--64~t TNT equivalent, for which the measured
acoustic energy would represent roughly 0.8--1.1\%. This is a sensitivity
comparison rather than an independently measured Falcon~9 yield.

Array back azimuths independently support SLC-40 as the source. Event
pressures scaled nearly proportionally with vector PGV, giving a median
seismic/acoustic coupling coefficient of
$2.3\times10^{-6}$~(m~s$^{-1}$)~Pa$^{-1}$, while H1 analysis identified
repeatable coupling bands near 4--5 and 20--24~Hz. No coherent regional
counterpart was confirmed. Of 171 station--sensor targets, 165 yielded usable
waveforms and 93 candidate phase windows at 67 stations were manually
reviewed; none produced coherent regional moveout. Contemporaneous upper-air
observations indicate strongly direction-dependent propagation, but detailed
regional ray paths were not modelled.

These observations show the value of close-range, continuously recording
seismo-acoustic instrumentation for reconstructing complex launch-vehicle
failures when engineering telemetry and direct site access are limited. A
permanent distributed KSC seismo-acoustic network would improve source
triangulation, constrain propagation and blast-wave decay, and provide repeat
measurements of local seismic/acoustic coupling. The archived waveforms,
catalogue, and reproducible workflow provide a reference dataset for that
monitoring design and for future studies of complex atmospheric explosions.

\begin{acknowledgements}
We thank Michael Wagner and USLaunchReport for recording the accident and for
written permission to use selected material in presentations and publications.
We thank Kennedy Space Center personnel for facilitating the BCHH deployment
and access to meteorological observations. Regional waveform data were
accessed through FDSN services; the IU Global Seismographic Network is cited
in the reference list \citep{asl_gsn_iu_1988}.
\end{acknowledgements}

\section*{Funding}

The present reanalysis and manuscript received no dedicated funding. SpaceX
provided US\$5,000 to USF in September 2016 for the original forensic analysis
and report. A contemporaneous US\$25,000 Florida Space Grant Consortium award
supported the broader research program, comprising US\$12,500 from the
Consortium and US\$12,500 in USF matching funds. Neither funder influenced the
present reanalysis, interpretation, manuscript, or decision to submit.

\section*{Data and code availability}

The Python notebooks, project modules, shared configuration, environment
specification, data manifest, output conventions, and derived products will be
archived as a versioned release in a repository issuing a persistent DOI. BCHH
waveforms, metadata, calibration information, manual picks, meteorological
inputs, event catalogue, and derived tables will be deposited in the same or a
linked record. The 48-h six-channel input used for the precursor screen is
separate from the short accident-analysis record and will be identified with
its exact coverage and processing provenance. The DOI and repository URL will be inserted before submission.

Regional waveforms remain available from their FDSN data centres and network
operators. The archive will include the station/channel manifest, availability
audit, screening metrics, and manual-review status/decisions. Selected regional
waveforms used for candidate validation will be retained as MiniSEED with
event-epoch StationXML and exact request/coverage metadata. The November~2016
technical report is cited as an unpublished historical source and is not
publicly available \citep{thompson_mcnutt_fireball_2016}; the 2017 AGU
presentation abstract is publicly archived \citep{thompson_agu_falcon9_2017}. The USLaunchReport and
Scott Manley videos remain at their cited URLs
\citep{uslaunchreport_amos6_video_2016,manley_spacex_explosion_2016}. Copyright
in the source footage remains with USLaunchReport.com/Michael Wagner; the
archive will provide synchronization metadata and code rather than redistribute
the video.

\section*{Competing interests}

The authors declare no competing interests. SpaceX support was limited to the
original 2016 forensic work; SpaceX had no role in the present reanalysis,
interpretation, manuscript preparation, or submission decision.

% The Seismica class selects abbrvnat_seismica_upcasetitle automatically.
\bibliography{mybibfile52}

\end{document}
